\documentclass[11pt,a4paper]{article}

\usepackage[utf8]{inputenc}
\usepackage[T1]{fontenc}
\usepackage[ngerman]{babel}
\usepackage{amsmath,amssymb,amsthm}
\usepackage[a4paper,margin=2.5cm]{geometry}

\title{EABC-Operatoralgebra und Renormierungskette\\[0.3em]
\large Ikosaeder--Dodekaeder-Dualität}
\author{Thomas Hoffbauer}
\date{}

\newtheorem{theorem}{Satz}
\newtheorem{definition}{Definition}
\newtheorem{remark}{Bemerkung}
\newtheorem{proposition}{Proposition}
\newtheorem{lemma}{Lemma}
\newtheorem{corollary}{Korollar}
\newtheorem{hypothesis}{Hypothese}
\newtheorem{conjecture}{Vermutung}

\newcommand{\R}{\mathbb{R}}
\newcommand{\Q}{\mathbb{Q}}
\newcommand{\Z}{\mathbb{Z}}

\begin{document}
\maketitle

\begin{abstract}
Eine EABC-zulässige Primzahl $p$ erzeugt auf dem Konfigurationsraum einen
messbaren Rang-1-Anisotropie-Defekt der Größe $w_p$:
\[
  p \;\longmapsto\; w_p \;=\; \Delta\!\bigl(24I_3 + w_p\,vv^{\mathsf T}\bigr),
  \qquad \|v\|=1.
\]
Die robuste Retraktion $R^*_{\mathrm{EABC}}$ führt gestörte Konfigurationen
zurück in den Fixpunkt-Teilraum $\mathcal{F}$; der effektive Tensor wird dabei
isotrop: $M_{\mathrm{eff}}(R^*(K^+))=24I_3$.
Der inhaltliche geometrische Kern
$M_{\mathrm{geom}}(I_{12},\sigma_{\mathrm{geo}})=24I_3$
ist von der projektiven Lesart des effektiven Tensors strikt getrennt und in
Lean~4 maschinell verifiziert (\texttt{sorry}\,$=\,0$).
\end{abstract}

%=============================================================================
\section{Einleitung: historische und moderne Einordnung}
\label{sec:introduction}
%=============================================================================

Die EABC-Renormierungskette verbindet diskrete Arithmetik, polyedrische Symmetrie
und effektive Tensorphysik in einem einzigen, formal prüfbaren Modell.
Mehrere historische Linien konvergieren hier --- ohne die Resultate dieser Arbeit
als deren unmittelbare Fortsetzung zu beanspruchen.

\paragraph{Leonhard Euler (1707--1783).}
Eulers Polyederformel und seine systematische Theorie der platonischen Körper
verankern das Ikosaeder als \emph{mathematisches} Referenzobjekt, nicht bloß als
Zeichnung.
Die zwölf Ecken des regulären Ikosaeders auf $S^2$ und die zugehörige
$A_5$-Symmetrie sind der geometrische Boden von $M_{\mathrm{geom}}$;
Theorem~\ref{thm:geom-fixpoint} ist in diesem Sinne ein Euler-artiger Satz:
eine endliche, explizite Summe symmetrischer dyadischer Produkte kollabiert
zu einem skalaren Vielfachen von $I_3$.
Die Peano-Schalenstapel (Abschnitt~\ref{sec:formal-verification}) erinnern an
Eulers Methode, diskrete Strukturen schichtweise über $\mathbb{N}$ zu erzeugen.

\paragraph{Bernhard Riemann (1826--1866).}
Riemanns Trennung von lokalen Koordinaten und globaler Struktur auf Mannigfaltigkeiten
entspricht methodisch der hier durchgeführten Separation von
$M_{\mathrm{eff}}$ (Konfigurationsraum) und $M_{\mathrm{geom}}$ (geometrischer Fixpunkt,
Abschnitt~\ref{sec:two-levels}).
Der Ordnungsparameter $\Delta(M)=\lambda_{\max}-\lambda_{\min}$ ist ein
spektral-geometrisches Funktional im riemannschen Sinne: Isotropie bedeutet
Entartung des Spektrums.
Die Primzahlen $p=2n+1$ treten als \emph{diskrete Normschalen-Indizes} auf ---
nicht als analytische Objekte der Riemannschen $\zeta$-Funktion.
Der strukturelle Bezug zu Riemann liegt in der Frage, wie globale Invarianz
($24I_3$) aus lokaler arithmetischer Störung ($w_p$) durch Projektion
wiederhergestellt wird.

\paragraph{Paul Dirac (1902--1984).}
Diracs Vakuum-Bild --- ein symmetrischer Grundzustand, dem diskrete
Anregungen (Massen) als Störungen überlagert sind --- spiegelt sich in der
Zwei-Ebenen-Architektur wider:
$M_{\mathrm{geom}}=24I_3$ ist das isotrope ``Vakuum'',
$M_{\mathrm{eff}}(K^+)=24I_3+w_pvv^{\mathsf T}$ die messbare Anisotropie
(Theorem~\ref{thm:prime-delta}, \texttt{rankOne\_defect\_not\_isotropic}).
Die Rang-1-Störung $w_pvv^{\mathsf T}$ ist die minimalste nicht-triviale
Symmetriebrechung in $\mathrm{Sym}_3(\mathbb{R})$; $\Delta(M^+)=w_p$ quantisiert
das ``Massen-Residuum'' der Störung.
Dirac würde die Retraktion $R^*$ nicht als Dynamik, sondern als
Zustandsreduktion lesen: Rückkehr zum symmetrischen Referenzzustand.

\paragraph{Gerard 't~Hooft (*1946).}
Die Renormierungs- und Eichtheorie-Tradition, deren Höhepunkt 't~Hoofts Arbeit
zur renormierbaren nicht-abelschen Eichfeldtheorie bildet, liefert das
konzeptuelle Vokabular des \emph{Retraktionsoperators} $R^*_{\mathrm{EABC}}$:
kein ad-hoc-Glätten, sondern projektive Rückführung auf einen wohldefinierten
Fixpunkt-Teilraum $\mathcal{F}$.
Die Idempotenz $(R^*)^2=R^*$ (Lemma~\ref{lem:idempotence}) entspricht dem
Gedanken, dass ein korrekt gewählter ``Gauge-Fix'' nach einem Schritt stabil bleibt.
Die EABC-Theorie beansprucht keine QFT-Renormierungsgruppe im analytischen Sinn;
sie übernimmt die \emph{Struktur}: lokaler Defekt, globaler Fixpunkt, ein Schritt genügt.

\paragraph{Zeitgenössische Nachbarn.}
Maynard, Tao und Scholze behandeln Primzahlen auf anderem Niveau
(Gaps, additive Kombinatorik, perfektoide Geometrie; Abschnitt~\ref{sec:modern-neighbors}).
Die EABC-Theorie bleibt diskret-modular: $p\bmod 12$ klassifiziert Normschalen,
nicht die Primzahlverteilung selbst.

%=============================================================================
\section{EABC-Familien und Anisotropietensor}
%=============================================================================

Die vier invertierbaren Restklassen modulo~$12$ definieren die EABC-Familien
\[
  E \equiv 1,\qquad A \equiv 5,\qquad B \equiv 7,\qquad C \equiv 11 \pmod{12},
\]
mit Gewichten $w_E=1$, $w_A=5$, $w_B=7$, $w_C=11$.

Sei $\{v_i\}_{i=1}^{12}\subset S^2\subset\R^3$ ein 12-Kuss-System (Einheitsrichtungen)
und $\sigma\colon\{1,\ldots,12\}\to\{E,A,B,C\}$ eine Label-Zuweisung.
Der \emph{Anisotropietensor} ist
\[
  M(\sigma)
  \;=\;
  \sum_{i=1}^{12} w_{\sigma(i)}\, v_i v_i^{\mathsf T}
  \;\in\;
  \mathrm{Sym}_3(\R).
\]

\begin{definition}[Ordnungsparameter]
\[
  \Delta(M) \;:=\; \lambda_{\max}(M) - \lambda_{\min}(M),
\]
wobei $\lambda_1\le\lambda_2\le\lambda_3$ die Eigenwerte von $M$ bezeichnen.
Isotropie entspricht $\Delta(M)=0$.
\end{definition}

%=============================================================================
\section{Renormierungsoperatoren}
%=============================================================================

Sei $K_n=(V_n,\sigma_n)$ eine Konfiguration aus Vertices $V_n$ und Zuweisung $\sigma_n$.

\begin{definition}[$R_{\mathrm{EABC}}$]
Der Operator $R_{\mathrm{EABC}}$ projiziert $\sigma$ auf die geometrisch zulässige
3+3+3+3-Zuweisung: pro Ecke $v_i$ wird die mod-$12$-Restklasse
$r(v_i)$ aus den skalierten Koordinaten und die nächste EABC-Restklasse
bestimmt. $R_{\mathrm{EABC}}$ ist idempotent; Fixpunkte erfordern perfekte
geometrische Konsistenz der Koordinaten-Klassifikation.
\end{definition}

\begin{definition}[$P_I$ (Voronoi-Pullback)]
$P_I$ bildet gestörte Kussrichtungen auf das Referenz-Ikosaeder ab:
\[
  P_I(v) \;=\; \arg\max_{r_j\in I_{12}} \langle v, r_j\rangle,
\]
wobei $I_{12}$ die 12 Referenz-Ecken des regulären Ikosaeders auf $S^2$ ist.
\end{definition}

\begin{definition}[$R^*_{\mathrm{EABC}}$]
Der robuste Fixpunktoperator ist die Komposition
\[
  R^*_{\mathrm{EABC}} \;:=\; P_I \circ R_{\mathrm{EABC}}.
\]
Zuerst geometrische Rückführung, dann Label-Projektion.
\end{definition}

\subsection{Konfigurationsraum und Retraktion}

\begin{definition}[Konfigurationsraum $\mathcal{C}$]
$\mathcal{C}$ ist die Menge aller Konfigurationen
\[
  K=(V,\sigma),
  \qquad
  V=(v_1,\ldots,v_{12})\subset S^2,
  \qquad
  \sigma\colon\{1,\ldots,12\}\to\{E,A,B,C\},
\]
also aller (ggf.\ gestörten) EABC-Kuss-Konfigurationen.
\end{definition}

\begin{definition}[Fixpunkt-Teilraum $\mathcal{F}$]
\[
  \mathcal{F}
  \;:=\;
  \bigl\{K\in\mathcal{C} : R^*_{\mathrm{EABC}}(K)=K\bigr\}
  \;\subset\;
  \mathcal{C}
\]
ist der Teilraum der \emph{robusten Fixpunkte} des Operators $R^*_{\mathrm{EABC}}$.
Das geometrische Referenzobjekt $I_{12}^{\mathrm{EABC}}$ liegt in $\mathcal{F}$.
\end{definition}

\begin{definition}[Einbettung und Retraktion]
Die kanonische Einbettung ist
\[
  \iota\colon \mathcal{F}\hookrightarrow \mathcal{C},\qquad \iota(K)=K.
\]
Der robuste Operator wirkt als Abbildung
\[
  R^*_{\mathrm{EABC}}\colon \mathcal{C}\longrightarrow \mathcal{F}.
\]
\end{definition}

\begin{proposition}[Retraktion auf $\mathcal{F}$]
\label{prop:retraction}
Mit $\iota$ und $R^*_{\mathrm{EABC}}$ wie oben gilt:
\begin{enumerate}
\item[(a)] \textbf{Idempotenz auf $\mathcal{C}$:}
  $(R^*_{\mathrm{EABC}})^2 = R^*_{\mathrm{EABC}}$.
\item[(b)] \textbf{Fixpunkt auf $\mathcal{F}$:}
  $R^*_{\mathrm{EABC}}(K)=K$ für alle $K\in\mathcal{F}$.
\item[(c)] \textbf{Retraktionsidentität:}
  $R^*_{\mathrm{EABC}}\circ \iota = \mathrm{id}_{\mathcal{F}}$.
\end{enumerate}
Insbesondere ist $R^*_{\mathrm{EABC}}\colon\mathcal{C}\to\mathcal{F}$ eine
\emph{Retraktion} von $\mathcal{C}$ auf den Fixpunkt-Teilraum $\mathcal{F}$.
\end{proposition}

\begin{proof}[Skizze]
(a) ist Lemma~\ref{lem:idempotence} unten.
Für (b) folgt aus der Definition von $\mathcal{F}$.
Für (c) setze $K\in\mathcal{F}$; dann ist
$R^*_{\mathrm{EABC}}(\iota(K))=R^*_{\mathrm{EABC}}(K)=K=\iota(K)$.
\end{proof}

\begin{proposition}[Konvergenz in einem Schritt]
\label{prop:convergence}
Ist $R^*_{\mathrm{EABC}}$ idempotent, so gilt für alle $K\in\mathcal{C}$ und
alle $n\ge 1$:
\[
  (R^*_{\mathrm{EABC}})^n(K) \;=\; R^*_{\mathrm{EABC}}(K).
\]
Insbesondere konvergiert die Renormierungsiteration nach \emph{einem} Schritt
auf das Bild $\mathcal{F}$.
\end{proposition}

\begin{proof}
Vollständig formalisiert in \texttt{eabc-renorm/EabcRenorm/Retraction.lean}
als \texttt{IsIdempotent.iterate\_ge\_one}.
Für $n=1$ trivial; für $n\ge 2$ folgt $(R^*)^n=(R^*)^{n-1}\circ R^*$
und Induktion über die Idempotenz $(R^*)^2=R^*$.
\end{proof}

%=============================================================================
\section{Fixpunkt und Defektklassen}
%=============================================================================

\begin{definition}[Geometrischer Fixpunkt $I_{12}^{\mathrm{EABC}}$]
Die Konfiguration $I_{12}^{\mathrm{EABC}}$ besteht aus den 12 Ikosaeder-Ecken
mit der geometrischen EABC-Zuweisung (je 3 Labels pro Familie).
\end{definition}

\begin{lemma}[Lemma 1A — Spur am balancierten Fixpunkt]
\label{lem:trace}
Sei $M(\sigma)=\sum_{i=1}^{12} w_{\sigma(i)}\, v_i v_i^{\mathsf T}$ mit $\|v_i\|=1$
und 3+3+3+3-Balance der Label-Zuweisung $\sigma$.
Dann gilt $\operatorname{tr}(M)=72$.
\end{lemma}

\begin{proof}
Da $\operatorname{tr}(v_i v_i^{\mathsf T})=\|v_i\|^2=1$, folgt
\[
  \operatorname{tr}(M)
  \;=\;
  \sum_{i=1}^{12} w_{\sigma(i)}
  \;=\;
  3(1+5+7+11)
  \;=\;
  72.
\]
\end{proof}

\begin{lemma}[Lemma 1B — Skalar aus Isotropie]
\label{lem:scalar}
Ist $M\in\mathrm{Sym}_3(\R)$ isotrop, d.\,h.\ $M=\lambda I$ für ein $\lambda\in\R$,
und gilt $\operatorname{tr}(M)=72$, dann folgt $\lambda=24$ und
$\lambda_1=\lambda_2=\lambda_3=24$.
\end{lemma}

\begin{proof}
Aus $M=\lambda I$ folgt $\operatorname{tr}(M)=3\lambda=72$, also $\lambda=24$.
\end{proof}

\begin{lemma}[Lemma 1C — Geometrische Zuordnung der mod-12-Bahnen]
\label{lem:orbit-assignment}
Die Ikosaeder-Ecken zerfallen bei Skala $s=6$ in vier Bahnen
$\mathcal{O}_2,\mathcal{O}_4,\mathcal{O}_8,\mathcal{O}_{10}$
(je drei Ecken).
Die Abbildung \emph{nächste EABC-Restklasse} liefert
\[
  \mathcal{O}_2\mapsto E\;(w=1),\quad
  \mathcal{O}_4\mapsto A\;(w=5),\quad
  \mathcal{O}_8\mapsto B\;(w=7),\quad
  \mathcal{O}_{10}\mapsto C\;(w=11).
\]
Diese Zuordnung ist injektiv auf den Familien und fest durch
$r(v)=(\mathrm{round}(sx)+\mathrm{round}(sy)+\mathrm{round}(sz))\bmod 12$
und die Kreis-Distanz auf $\Z/12\Z$.
\end{lemma}

\begin{proof}
Die vier Aussagen $\mathcal{O}_r\mapsto\{\text{E,A,B,C}\}$ sind endliche
Verifikationen auf den Restklassen $2,4,8,10$; formalisiert in
\texttt{eabc-renorm/EabcRenorm/Mod12.lean} per \texttt{native\_decide}.
\end{proof}

\begin{lemma}[Lemma 1D — Geometrische Isotropie]
\label{lem:geom-iso}
Am geometrischen Fixpunkt $I_{12}^{\mathrm{EABC}}$ gilt
\[
  M(\sigma)
  \;=\;
  \sum_{i=1}^{12} w_{\sigma(i)}\, v_i v_i^{\mathsf T}
  \;=\;
  24\,I_3.
\]
\end{lemma}

\begin{proof}[Zwei unabhängige Wege]
\textbf{Weg I (direkte Summe).}
Mit den expliziten Koordinaten aus Anhang~\ref{app:generators} ergibt die
Summe der zwölf gewichteten äußeren Produkte in $\Q(\sqrt5)$ exakt $24I_3$.

\smallskip\noindent
\textbf{Weg II (Symmetrie + Schur).}
Sei $\mathcal{I}\cong A_5$ die Ikosaeder-Rotationsgruppe.
Die gewichtete Bahn-Summe $M=\sum_r w_{\phi(r)}T_r$ ist $\mathcal{I}$-invariant;
nach Schur folgt $M=\lambda I$.
Lemma~\ref{lem:trace} liefert $\lambda=24$.
Die Erzeuger-Verifikation steht in Anhang~\ref{app:generators}.
\end{proof}

\begin{theorem}[Isotropie-Lemma — Komposition]
\label{thm:isotropy}
Am geometrischen EABC-Fixpunkt $I_{12}^{\mathrm{EABC}}$ hat $M(\sigma)$ die
Eigenwerte $\lambda_1=\lambda_2=\lambda_3=24$, also $\Delta(M)=0$.
\end{theorem}

\begin{proof}
Lemma~\ref{lem:geom-iso} liefert $M=24I_3$ direkt.
Alternativ: Lemma~\ref{lem:geom-iso} (Weg II) gibt Isotropie,
Lemma~\ref{lem:trace} gibt $\operatorname{tr}(M)=72$,
Lemma~\ref{lem:scalar} gibt $\lambda=24$.
\end{proof}

\begin{remark}[Lean-Formalisierung (V4)]
In \texttt{eabc-renorm} ist der Beweis dreistufig angelegt:
\texttt{trace\_geometric\_fixpoint} (1A),
\texttt{isotropy\_from\_trace} (1B),
\texttt{geometric\_fixpoint\_isotropy} (1C/1D).
Stand: \texttt{lake build} grün, \texttt{sorry}\,$=\,0$
(siehe Abschnitt~\ref{sec:formal-verification}).
\end{remark}

\begin{remark}[Numerische Bestätigung]
Unabhängig bestätigt \texttt{eabc\_icosahedron\_test.py},
\texttt{run\_core\_lemma\_tests()} die Aussage (4/4 PASS, Toleranz $10^{-6}$).
\end{remark}

\begin{lemma}[Projektor-Idempotenz]
\label{lem:idempotence}
$R^*_{\mathrm{EABC}} = P_I \circ R_{\mathrm{EABC}}$ ist idempotent:
\[
  (R^*_{\mathrm{EABC}})^2 \;=\; R^*_{\mathrm{EABC}}.
\]
\end{lemma}

\begin{proof}[Skizze]
Nach $P_I$ liegen alle Vertices auf den Referenz-Ecken; $R_{\mathrm{EABC}}$ ist
auf dieser Untervarietät bereits Fixpunktabbildung. Ein zweiter Schritt ändert
weder Geometrie noch Labels.
\end{proof}

\smallskip\noindent
\textit{Numerisch verifiziert in \texttt{eabc\_icosahedron\_test.py},
\texttt{run\_core\_lemma\_tests()} (4/4 PASS, Toleranz $10^{-6}$).}

\begin{definition}[Defektklassen]
\begin{align*}
  \text{FIXPOINT}            &\colon \Delta(M)\approx 0,\; 3{+}3{+}3{+}3\text{-Balance},\; \text{ungestörte Geometrie},\\
  \text{LABEL\_DEFECT}       &\colon \text{Label-Zuweisung ohne 3+3+3+3-Balance (vor } R_{\mathrm{EABC}}\text{)},\\
  \text{DIRECTION\_DEFECT}   &\colon \text{balancierte Labels, aber }\Delta(M)>\varepsilon \text{ bei gestörter Geometrie},\\
  \text{CLASSIFICATION\_DEFECT} &\colon \text{mod-}12\text{-Balance der geometrischen Zuweisung gebrochen},\\
  \text{PRIME\_DEFECT}(p)    &\colon p=2n+1\in\mathbb{P},\; p\bmod 12\in\{1,5,7,11\},\;
                               \text{neue Prim-Normkugel verletzt zunächst die 3+3+3+3-Balance}.
\end{align*}
\end{definition}

%=============================================================================
\section{Prim-Erzeugungsschritt}
\label{sec:prime-step}
%=============================================================================

\subsection{Operation $2n\mapsto 2n+1$}

Sei $N=2n$ eine gerade Normschale.
Der \emph{Prim-Erzeugungsschritt} bildet
\[
  P(2n)
  \;:=\;
  \begin{cases}
    2n+1, & \text{falls }2n+1\in\mathbb{P},\\
    \bot, & \text{sonst}.
  \end{cases}
\]
Für $p=P(2n)$ gilt bei Primzahlen $p>3$ stets
$p\bmod 12\in\{1,5,7,11\}$; die neue Normkugel landet damit in genau
einer EABC-Familie.

\begin{definition}[Erweiterung um eine Prim-Normkugel]
Gegeben $K_n\in\mathcal{F}$ mit Label-Zuweisung $\sigma$ und $p=P(2n)$.
Die erweiterte Konfiguration $K_n^+$ entsteht durch Hinzufügen einer
Einheitsrichtung $v\in S^2$ mit Label $\mathrm{EABC}(p)$:
\[
  M^+ \;=\; M + w_p\, v v^{\mathsf T},
  \qquad
  w_p\in\{1,5,7,11\}\text{ gemäß }p\bmod 12.
\]
Die Familienbesetzung wird von $(3,3,3,3)$ zu $(4,3,3,3)$ (in genau einer Komponente).
\end{definition}

\begin{proposition}[PRIME\_DEFECT $\Rightarrow$ LABEL\_DEFECT]
\label{prop:prime-implies-label}
Ist $\text{PRIME\_DEFECT}(p)$ aktiv, so gilt $K_n^+\notin\mathcal{F}$ und
$R^*_{\mathrm{EABC}}(K_n^+)\neq K_n^+$.
Die Störung ist ein \emph{arithmetisch privilegierter} Label-Defekt.
\end{proposition}

\subsection{Anisotropie nach dem Primschritt}

Am Fixpunkt war $M=24I_3$.
Nach dem Primschritt gilt im Rang-1-Modell
\[
  M^+ \;=\; 24\,I_3 + w_p\, v v^{\mathsf T},
  \qquad \|v\|=1.
\]

\begin{theorem}[Primdefekt-Satz]
\label{thm:prime-delta}
Für $w_p>0$ und $\|v\|=1$ hat $M^+=24I_3+w_p vv^{\mathsf T}$ die Eigenwerte
$\lambda_1=\lambda_2=24$ und $\lambda_3=24+w_p$ (bzw.\ Permutation falls $v$ nicht
Koordinatenachse ist).
Insbesondere
\[
  \Delta(M^+) \;=\; w_p \;\in\; \{1,5,7,11\}.
\]
\end{theorem}

\begin{proof}
Für symmetrische Rang-1-Störungen $w\,vv^{\mathsf T}$ mit $\|v\|=1$ ist $v$ Eigenvektor
zum Eigenwert $w$; die Orthogonalen zu $v$ haben Eigenwert $0$.
Also hat $24I_3+w_p vv^{\mathsf T}$ die Eigenwerte $24$ (doppelt) und $24+w_p$ (einfach).
Damit $\Delta(M^+)=\lambda_{\max}-\lambda_{\min}=w_p$.
Formalisiert in \texttt{eabc-renorm/EabcRenorm/PrimeDefect.lean}
als \texttt{delta\_rank\_one\_perturbation}.
\end{proof}

\begin{remark}[Primzahlen als Symmetriebrecher]
Eine Primzahl erscheint hier nicht als bloße Zahl ohne Teiler, sondern als
\emph{neue ungerade Normkugel}, die aus einer geraden Schale hervorgeht,
ein zulässiges EABC-Label trägt, die isotrope Balance bricht und einen
Renormierungsimpuls auslöst.
\end{remark}

\subsection{$R^*_{\mathrm{EABC}}$ als Prim-Entstörungsoperator}

Bei einem Primschritt hat $R^*_{\mathrm{EABC}}=P_I\circ R_{\mathrm{EABC}}$ die Bedeutung
\[
  \text{Prim-Erzeugung}
  \;\Rightarrow\;
  \text{PRIME\_DEFECT}
  \;\Rightarrow\;
  \text{projektive Rückführung}.
\]
Die Kette lautet
\[
  K_n \;\longrightarrow\; K_n^+ \;\longrightarrow\; R^*_{\mathrm{EABC}}(K_n^+).
\]
Nach Proposition~\ref{prop:convergence} genügt \emph{ein} Projektionsschritt:
$(R^*_{\mathrm{EABC}})^2=R^*_{\mathrm{EABC}}$.

\begin{theorem}[Prim-Normkugel-Satz — Label-Teil]
\label{thm:prime-norm}
Sei $K_n\in\mathcal{F}$ mit geometrischer Label-Zuweisung und
$K_n^+=\mathrm{Extend}(p)\,K_n$ für $p=P(2n)\in\mathbb{P}$, $p>3$.
Dann gilt:
\begin{enumerate}
\item[(a)] $\text{PRIME\_DEFECT}(p)$ und $K_n^+\notin\mathcal{F}$;
\item[(b)] $R^*_{\mathrm{EABC}}(K_n^+)\in\mathcal{F}$ mit wiederhergestellter
  3+3+3+3-Balance;
\item[(c)] ein Schritt genügt: $(R^*_{\mathrm{EABC}})^n(K_n^+)=R^*_{\mathrm{EABC}}(K_n^+)$
  für alle $n\ge 1$.
\end{enumerate}
\end{theorem}

\begin{proof}[Beweisidee]
\textbf{(a)} Die virtuelle Prim-Normkugel erhöht die Multiketten-Zählung einer
Familie von $3$ auf $4$ bei unveränderter 12er-Geometrie
(Lemma~\ref{lem:prime-imbalance}, formalisiert als
\texttt{extendPrime\_creates\_defect}).

\smallskip\noindent
\textbf{(b)} Entscheidend ist, dass $R_{\mathrm{EABC}}$ \emph{labelblind} ist:
Es setzt die Zuweisung unabhängig von der Eingabe auf die geometrische
3+3+3+3-Struktur und löscht den Prim-Schalen-Defekt
(\texttt{rEabc\_restores\_balance} in \texttt{PrimeDefect.lean}).
Nach $P_I$ liegt die Geometrie wieder auf dem Referenz-Ikosaeder.

\smallskip\noindent
\textbf{(c)} Folgt aus der Idempotenz von $R^*_{\mathrm{EABC}}$
(Proposition~\ref{prop:convergence}).
\end{proof}

\subsection{Die zwei Tensorebenen: Geometrischer Fixpunkt versus effektiver Konfigurationsraum}
\label{sec:two-levels}

Ein zentrales Missverständnis bei diskreten Symmetriebrechungen liegt in der
unzulässigen Vermischung von physikalischem Zustand und geometrischem Hintergrund.
Das EABC-Renormierungsmodell löst diesen Konflikt durch eine strikte Separation
zweier Tensorebenen --- auf jeder Peano-Schale $n\in\mathbb{N}$ und zunächst
unabhängig von metrischen Radien (Phase~A in Abschnitt~\ref{sec:formal-verification}).

Der Primdefekt wirkt gleichzeitig auf Label- und Tensorebene:
\begin{enumerate}
\item[(L)] Label-Multikette: $(4,3,3,3)$ statt $(3,3,3,3)$;
\item[(T)] Tensor: $M_{\mathrm{eff}}(K^+)=24I_3+w_pvv^{\mathsf T}$ (Rang-1-Störung).
\end{enumerate}

\subsubsection{Die rein geometrische Ebene: der Fixpunkt $M_{\mathrm{geom}}$}

Die Basis bildet der rein geometrische Tensor $M_{\mathrm{geom}}$, der ausschließlich
die ungestörte Ikosaeder-Konfiguration abbildet.
Für Vertices $V=(v_i)_{i=1}^{12}$ und Label-Zuweisung $\sigma$ gilt
(Definition~\ref{def:m-geom}, \texttt{M\_geom} in \texttt{EffectiveTensor.lean})
\[
  M_{\mathrm{geom}}(V,\sigma)
  \;=\;
  \sum_{i=1}^{12} w_{\sigma(i)}\, v_i v_i^{\mathsf T}.
\]
Das in \texttt{IcosahedronSymmetry.lean} verifizierte Theorem
\texttt{geometric\_fixpoint\_isotropy} beweist, dass für das Referenzobjekt
$I_{12}^{\mathrm{EABC}}$ mit $\sigma_{\mathrm{geo}}$ alle Richtungsabhängigkeiten
exakt kollabieren:
\[
  \boxed{
  M_{\mathrm{geom}}(I_{12},\sigma_{\mathrm{geo}}) \;=\; 24I_3
  }.
\]
Dieser Zustand ist perfekt isotrop und invariant unter $SO(3)$.
Er enthält keine dynamischen Defektfelder; in Lean entspricht er
\texttt{geometricM} bzw.\ \texttt{M\_geom\_geometric\_fixpoint}.

\subsubsection{Die effektive physikalische Ebene: der Konfigurationsraum $M_{\mathrm{eff}}$}

Der effektive Tensor $M_{\mathrm{eff}}$ operiert auf dem Konfigurationsraum
$\mathcal{C}$ (Lean: \texttt{Config} mit Feldern \texttt{label},
\texttt{primeShell}, \texttt{tensorDefect}).
Tritt ein EABC-Primereignis $p$ mit Gewicht $w_p\in\{1,5,7,11\}$ entlang eines
normierten Defektvektors $v$ ($\|v\|=1$) auf, so setzt man
(Definition~\ref{def:m-eff})
\[
  K^+ \;:=\; \mathrm{Extend}(p,v)\,K,
  \qquad
  M_{\mathrm{eff}}(K^+) \;=\; 24I_3 + w_p\,vv^{\mathsf T}.
\]
In Lean: \texttt{extendPrimeTensor} schreibt \texttt{tensorDefect = some (w\_p, v)},
und \texttt{M\_eff\_with\_defect} liefert \texttt{rankOnePerturbation 24 w v}.
Das Theorem \texttt{rankOne\_defect\_not\_isotropic} zeigt: für $w_p>0$ ist
$M_{\mathrm{eff}}(K^+)$ nicht isotrop.
Die messbare Anisotropie ist
\[
  \Delta\bigl(M_{\mathrm{eff}}(K^+)\bigr) \;=\; w_p
\]
(Theorem~\ref{thm:prime-delta}, \texttt{delta\_prime\_defect\_at\_fixpoint}).

\subsubsection{Die Wirkung von $R^*_{\mathrm{EABC}}$: ein reiner Projektionssatz}

$R^*_{\mathrm{EABC}}$ ist kein Dämpfungs- oder Glättungsoperator, sondern ein
algebraischer Retrakt auf $\mathcal{F}$.
Die Tensor-Restauration geschieht nicht analytisch \emph{innerhalb} der gestörten
Matrix $M_{\mathrm{eff}}(K^+)$, sondern durch projektiven Kollaps der Konfiguration
zurück auf die geometrische Referenz:
\[
  K^+
  \;\xrightarrow{\;R^*_{\mathrm{EABC}}\;}
  (I_{12},\sigma_{\mathrm{geo}})
  \;\xrightarrow{\;M_{\mathrm{geom}}\;}
  24I_3.
\]
Formal faktorisiert die Kette in zwei unabhängige Schritte:
\begin{enumerate}
\item \textbf{Projektion (Ebene~1):}
  $M_{\mathrm{eff}}(R^*(K))=M_{\mathrm{geom}}(I_{12},\sigma_{\mathrm{geo}})$
  ist \emph{definitorisch}
  (Lemma~\ref{lem:tensor-projection}, \texttt{M\_eff\_projection}:
  \texttt{rStar} setzt \texttt{tensorDefect} auf \texttt{none}).
\item \textbf{Fixpunkt (Ebene~2):}
  $M_{\mathrm{geom}}(I_{12},\sigma_{\mathrm{geo}})=24I_3$
  ist der \emph{inhaltliche} Ikosaeder-Satz
  (Theorem~\ref{thm:geom-fixpoint}).
\end{enumerate}
Für eine einzelne Schale liefert \texttt{prime\_norm\_full\_restoration} beides:
$R^*(K^+)\in\mathcal{F}$ und $M_{\mathrm{eff}}(R^*(K^+))=24I_3$.
Auf einem Peano-Schalenstapel gilt dasselbe komponentenweise
(\texttt{all\_shells\_tensor\_restored} in \texttt{ShellStack.lean}).

\begin{table}[h!]
\centering
\small
\begin{tabular}{p{0.14\linewidth}p{0.36\linewidth}p{0.22\linewidth}p{0.22\linewidth}}
\hline
\textbf{Ebene} & \textbf{Mathematisches Objekt} & \textbf{Symmetrie} & \textbf{Lean} \\
\hline
Geometrisch &
  $M_{\mathrm{geom}}(I_{12},\sigma_{\mathrm{geo}})=24I_3$ &
  Isotropie (Referenz) &
  \texttt{geometric\_fixpoint\_isotropy} \\
Physikalisch &
  $M_{\mathrm{eff}}(K^+)=24I_3+w_pvv^{\mathsf T}$ &
  Rang-1-Anisotropie &
  \texttt{extendPrimeTensor}, \texttt{rankOne\_defect\_not\_isotropic} \\
Projektiv &
  $M_{\mathrm{eff}}(R^*(K^+))=24I_3$ &
  restaurierte Isotropie &
  \texttt{prime\_norm\_full\_restoration} \\
Peano (global) &
  $\mathcal{R}^*(\mathcal{S}_{\mathrm{def}})=\mathcal{S}_{\mathrm{iso}}$ &
  schalenweise unabhängig &
  \texttt{all\_shells\_tensor\_restored} \\
\hline
\end{tabular}
\caption{Gegenüberstellung geometrischer Invarianz, effektiver Störung und projektiver Restauration.
\texttt{IsEabcFixpoint} beschreibt den Fixpunkt $K\in\mathcal{F}$, nicht direkt den Tensor;
die Isotropie $24I_3$ folgt über \texttt{M\_eff} und Theorem~\ref{thm:geom-fixpoint}.}
\label{tab:two-levels}
\end{table}

\subsection{Tensor-Restauration als Projektionssatz}
\label{sec:tensor-projection}

Der Mechanismus ist keine dynamische Tensor-Glättung, sondern die Faktorisierung
\[
  \mathcal{C}
  \;\xrightarrow{\;R^*\;}
  \mathcal{F}
  \;\xrightarrow{\;M_{\mathrm{geom}}\;}
  24I_3,
  \qquad
  K
  \;\stackrel{R^*}{\longmapsto}\;
  (I_{12},\sigma_{\mathrm{geo}})
  \;\stackrel{M_{\mathrm{geom}}}{\longmapsto}\;
  24I_3.
\]

\begin{definition}[Effektiver Tensor $M_{\mathrm{eff}}$]
\label{def:m-eff}
Auf dem Konfigurationsraum $\mathcal{C}$ setzt man
\[
  M_{\mathrm{eff}}(K)
  \;:=\;
  \begin{cases}
    M_{\mathrm{geom}}(I_{12},\sigma_{\mathrm{geo}}), & \text{kein Tensordefekt},\\[4pt]
    24I_3+w_pvv^{\mathsf T}, & \text{Primdefekt mit }w_p,\ \|v\|=1.
  \end{cases}
\]
(\texttt{M\_eff} in \texttt{EffectiveTensor.lean}).
\end{definition}

\begin{definition}[Geometrischer Tensor $M_{\mathrm{geom}}$]
\label{def:m-geom}
Für Vertices $V=(v_i)_{i=1}^{12}$ und Label-Zuweisung $\sigma$:
\[
  M_{\mathrm{geom}}(V,\sigma)
  \;:=\;
  \sum_{i=1}^{12} w_{\sigma(i)}\, v_i v_i^{\mathsf T}.
\]
\end{definition}

\begin{lemma}[Projektionssatz — effektive Ebene]
\label{lem:tensor-projection}
$R^*_{\mathrm{EABC}}$ löscht den Tensordefekt per Definition der Konfiguration:
\[
  M_{\mathrm{eff}}\bigl(R^*_{\mathrm{EABC}}(K)\bigr)
  \;=\;
  M_{\mathrm{geom}}(I_{12},\sigma_{\mathrm{geo}}).
\]
(\texttt{M\_eff\_projection}).
\end{lemma}

\begin{theorem}[Ikosaeder-Fixpunkt — geometrische Ebene]
\label{thm:geom-fixpoint}
Am geometrischen Referenz-Fixpunkt gilt
\[
  \boxed{
  M_{\mathrm{geom}}(I_{12},\sigma_{\mathrm{geo}})
  \;=\;
  24I_3
  }.
\]
(\texttt{geometric\_fixpoint\_isotropy}; Beweis via $R_3$-Invarianz,
Symmetrie, $\mathrm{tr}(M)=72$ und $M_{01}=0$ aus der gewichteten Summe der
vier $(\pm1,\pm\varphi,0)$-Ecken).
\end{theorem}

\begin{corollary}[Tensor-Restauration nach Primdefekt]
\label{cor:tensor-restore}
Für $K_n\in\mathcal{F}$, $p=P(2n)$ und $\|v\|=1$ gilt die Satzkette
\begin{align*}
  M_{\mathrm{eff}}(K_n^+) &= 24I_3+w_pvv^{\mathsf T},
  & \Delta(M_{\mathrm{eff}}(K_n^+)) &= w_p,
  \\
  w_p>0 &\Rightarrow M_{\mathrm{eff}}(K_n^+)\ \text{nicht isotrop},
  & R^*_{\mathrm{EABC}}(K_n^+) &\in \mathcal{F},
  \\
  M_{\mathrm{geom}}(I_{12},\sigma_{\mathrm{geo}}) &= 24I_3,
  &
  M_{\mathrm{eff}}(R^*_{\mathrm{EABC}}(K_n^+)) &= 24I_3.
\end{align*}
\end{corollary}

\begin{proof}
Theorem~\ref{thm:prime-delta} und \texttt{delta\_prime\_defect\_at\_fixpoint}
(vor $R^*$),
\texttt{rankOne\_defect\_not\_isotropic} (Nicht-Isotropie),
Theorem~\ref{thm:prime-norm} bzw.\ \texttt{prime\_norm\_full\_restoration}
(Label und Tensor nach $R^*$),
Lemma~\ref{lem:tensor-projection} (definitorisch, Ebene~1),
Theorem~\ref{thm:geom-fixpoint} (geometrisch, Ebene~2).
\end{proof}

\begin{remark}[Gutachter-Einordnung]
$M_{\mathrm{eff}}(R^*(K))=24I_3$ ist \emph{kein} mystischer Spektralsatz:
die Projektion ersetzt die defekte Konfiguration durch den geometrischen Fixpunkt;
dessen Isotropie ist der inhaltliche Satz (Theorem~\ref{thm:geom-fixpoint}).
$P_I$ (Voronoi-Pullback gestörter Vertices) ist eine geometrische Realisierung
von $R^*$, nicht der konzeptionelle Kern.
\end{remark}

\begin{remark}[$P_I$ als geometrische Realisierung (optional)]
\label{rem:pi-optional}
Für gestörte Richtungen $v\in S^2$ setzt
$P_I(v):=\arg\max_{r_j\in I_{12}}\langle v,r_j\rangle$.
Bei hinreichend kleiner Kuss-Störung zieht $P_I$ die Vertex-Menge auf $I_{12}$
zurück; zusammen mit labelblindem $R_{\mathrm{EABC}}$ erhält man dieselbe
Projektionssemantik wie oben.
\end{remark}

\begin{lemma}[Multiketten-Imbalance nach dem Primschritt]
\label{lem:prime-imbalance}
Aus $K_n\in\mathcal{F}$ und $p=P(2n)$ folgt für die Familienzählung mit virtueller
Primkugel $(4,3,3,3)$ statt $(3,3,3,3)$.
\end{lemma}

%=============================================================================
\section{Formale Verifikation}
\label{sec:formal-verification}
%=============================================================================

Die zentrale Satzkette ist in \texttt{eabc-renorm} (Lean~4.29, Mathlib~4.29)
maschinell verifiziert; \texttt{lake build} ist grün, \texttt{sorry}\,$=\,0$.

\begin{table}[h!]
\centering
\small
\begin{tabular}{p{0.38\linewidth}p{0.30\linewidth}p{0.24\linewidth}}
\hline
\textbf{Mathematische Aussage} & \textbf{Lean-Theorem} & \textbf{Datei} \\
\hline
$M_{\mathrm{geom}}(I_{12},\sigma_{\mathrm{geo}})=24I_3$
  & \texttt{geometric\_fixpoint\_isotropy} & \texttt{IcosahedronSymmetry.lean} \\
$M^+=24I_3+w_pvv^{\mathsf T}$
  & \texttt{rankOnePerturbation} & \texttt{PrimeDefect.lean} \\
$\Delta(M^+)=w_p$
  & \texttt{delta\_prime\_defect\_at\_fixpoint} & \texttt{PrimeDefect.lean} \\
$w_p>0\Rightarrow M^+$ nicht isotrop
  & \texttt{rankOne\_defect\_not\_isotropic} & \texttt{TensorRestoration.lean} \\
$R^*(K^+)\in\mathcal{F}$, $M_{\mathrm{eff}}=24I_3$
  & \texttt{prime\_norm\_full\_restoration} & \texttt{TensorRestoration.lean} \\
$M_{\mathrm{eff}}(R^*(K))=M_{\mathrm{geom}}(\mathrm{ref})$
  & \texttt{M\_eff\_projection} & \texttt{EffectiveTensor.lean} \\
\hline
\multicolumn{3}{l}{\textbf{Peano-Erweiterung (Multi-Schale, Phase~A)}} \\
\hline
Schalenstapel $\mathcal{S}$
  & \texttt{structure ShellStack} & \texttt{ShellStack.lean} \\
Lokaler Schalen-Retrakt
  & \texttt{shell\_defect\_restoration} & \texttt{ShellStack.lean} \\
Globaler Projektionssatz
  & \texttt{all\_shells\_restored} & \texttt{ShellStack.lean} \\
Globale Tensor-Restauration
  & \texttt{all\_shells\_tensor\_restored} & \texttt{ShellStack.lean} \\
\hline
\end{tabular}
\caption{Verifizierte Satzkette inkl.\ Peano-Schalenstapel.}
\label{tab:lean-theorems}
\end{table}

\subsection{Peano-Schachtelung und globaler Retrakt}

Eine einzelne Schale ist in Lean abgeschlossen.
Für geschachtelte Normschalen wird der Konfigurationsraum zu einem
\emph{indexbasierten Stapel} $\mathcal{S}=\{K^{(i)}\}_{i<n}$ über
$\mathbb{N}$ (Peano-Rekursion) erweitert.
Der globale Operator $\mathcal{R}^*$ wirkt komponentenweise:
auf markierten Schalen $i$ mit aktivem Primdefekt gilt
\[
  \mathcal{R}^*(K^{(i)}) \;=\; R^*_{\mathrm{EABC}}\!\bigl(\mathrm{Extend}(p)\,K^{(i)}\bigr),
\]
auf unberührten Schalen bleibt $K^{(i)}$ unverändert.
Da jede Schale den lokalen Satz \texttt{shell\_defect\_restoration} erfüllt,
folgt der globale Satz \texttt{all\_shells\_restored} ohne zusätzliche
Analyse-Hypothesen.

\subsection{Metrische Interpretation (Phase~B, optional)}

Die kombinatorische Verifikation hängt \emph{nicht} von metrischen Radien ab.
Für physikalische Lesarten lässt sich dennoch eine rein interpretative
Skalierungsfunktion einführen, z.\,B.
\[
  r_n \;=\; \phi^n\,r_0,
  \qquad \phi=\tfrac{1+\sqrt5}{2},
\]
ohne den algebraischen Restaurationssatz zu verändern:
\texttt{all\_shells\_restored} bleibt exakt, weil $R^*$ auf der
Label-/Defektebene operiert und $M_{\mathrm{geom}}=24I_3$ metrikunabhängig ist.

\subsection{Zeitgenössische Nachbarn: Maynard, Tao, Scholze}
\label{sec:modern-neighbors}

Die EABC-Theorie operiert nicht auf dem Niveau von Gaps zwischen Primzahlen
(Maynard--Tao) oder perfektoider Geometrie (Scholze), sondern auf einer
\emph{diskreten mod-$12$-Klassifikation} ungerader Normkugeln.
Der strukturelle Anschluss ist arithmetisch, nicht analytisch:
Primzahlen $p>3$ mit $p\bmod 12\in\{1,5,7,11\}$ erzeugen privilegierte
Defekt-Ereignisse mit quantisierten Gewichten $w_p\in\{1,5,7,11\}$.
Die Verifikation betrifft damit Symmetrie-Fixpunkte und projektive Retraktion,
nicht die Dichte oder Verteilung der Primzahlen selbst.
Im historischen Bogen (Abschnitt~\ref{sec:introduction}) steht diese Arbeit
näher bei Euler und 't~Hooft (diskrete Symmetrie, projektive Stabilisierung)
als bei den genannten Primzahlprogrammen der Gegenwart.

%=============================================================================
\section{Dualität}
%=============================================================================

Ikosaeder--Dodekaeder-Dualität: $I_{12}\leftrightarrow D_{20}$ (12 Ecken $\leftrightarrow$ 20 Flächen).
Unter Dualitätsabbildung und Label-Pushforward gilt:

\begin{lemma}[Dualitätsstabilität]
$\Delta(M)$ am geometrischen Fixpunkt ist unter Ikosaeder$\leftrightarrow$Dodekaeder-Dualität
invariant. Dagegen sind der Dipol $D=\sum_i w_{\sigma(i)} v_i$ und die Holonomie $H$
\emph{nicht} dualitätsinvariant.
\end{lemma}

\smallskip\noindent
\textit{Numerisch verifiziert in \texttt{eabc\_icosahedron\_test.py},
\texttt{run\_core\_lemma\_tests()} (4/4 PASS, Toleranz $10^{-6}$).}

%=============================================================================
\section{Renormierungskette}
%=============================================================================

Die diskrete Kette (ohne Primschritt) lautet
\[
  K_n \;\longrightarrow\; K_{n+1}
  \;=\;
  \underbrace{\text{Volumen-Renormierung}}_{\lambda_{\mathrm{sphere}}}
  \;\circ\;
  \underbrace{R^*_{\mathrm{EABC}}}_{P_I \circ R_{\mathrm{EABC}}}
  \;\circ\;
  \underbrace{D_{20}}_{\text{Dual}}
  \;\circ\;
  \underbrace{I_{12}}_{\text{Ikosaeder}}
  \;\circ\;
  \underbrace{\text{12-Kuss}}_{\text{Einheitsrichtungen}}.
\]

\paragraph{Erweiterung mit Prim-Erzeugungsschritt (Abschnitt~\ref{sec:prime-step}).}
Wenn $P(2n)=p$ prim, setze
\[
  K_n^+ \;:=\; \mathrm{Extend}(p)\,K_n,
  \qquad
  K_{n+1}
  \;=\;
  \lambda_{\mathrm{sphere}}
  \circ R^*_{\mathrm{EABC}}
  \circ \mathrm{Extend}(P)
  \;(K_n),
\]
wobei $\mathrm{Extend}(p)$ die Konfiguration um die Prim-Normkugel mit
Label $\mathrm{EABC}(p)$ erweitert.
Äquivalent:
\[
  K_n \;\longrightarrow\; K_n^+ \;\longrightarrow\;
  R^*_{\mathrm{EABC}}(K_n^+) \;\longrightarrow\; K_{n+1}.
\]

\begin{definition}[Volumen-Renormierung]
Für äußeren Radius $R>0$ setze $R_{\mathrm{out}}=R(1+2\pi)$ und
\[
  \lambda_{\mathrm{sphere}}
  \;=\;
  \left(\frac{V_0}{V_{\mathrm{out}}}\right)^{\!1/3}
  \;=\;
  \frac{1}{1+2\pi},
\]
mit $V_0=\tfrac{4}{3}\pi R^3$, $V_{\mathrm{out}}=\tfrac{4}{3}\pi R_{\mathrm{out}}^3$.
Die Vertex-Skalierung erfolgt isotrop: $v_i\mapsto \lambda_{\mathrm{sphere}}\,v_i$.
\end{definition}

\begin{lemma}[Volumen-Kompatibilität]
Isotrope Skalierung mit $\lambda_{\mathrm{sphere}} = 1/(1+2\pi)$ multipliziert alle
Eigenwerte von $M$ mit $\lambda_{\mathrm{sphere}}^2$; $\Delta(M)$ bleibt erhalten.
\end{lemma}

\smallskip\noindent
\textit{Numerisch verifiziert in \texttt{eabc\_icosahedron\_test.py},
\texttt{run\_core\_lemma\_tests()} (4/4 PASS, Toleranz $10^{-6}$).}

%=============================================================================
\section{Numerische Befunde}
%=============================================================================

Implementierung in \texttt{eabc\_icosahedron\_test.py}:

\begin{itemize}
\item $\varepsilon=0$: Nach $R_{\mathrm{EABC}}$ bzw.\ $R^*_{\mathrm{EABC}}$
  gilt $\Delta(M)\approx 0$ (Fixpunkt erreicht).
\item $\varepsilon>0$: Gestörte Geometrie ändert mod-$12$-Klassifikation;
  $R_{\mathrm{EABC}}$ allein reicht nicht --- gekoppelter Operator $R^*_{\mathrm{EABC}}$
  (Pullback + Label-Projektion) konvergiert zuverlässiger zu $\Delta(M)\approx 0$.
\item Gekoppelt vs.\ $R_{\mathrm{EABC}}$ allein: Bei $\varepsilon>0$ liefert
  $R^*_{\mathrm{EABC}}$ niedrigere mittlere $\Delta(M)$ und höhere Fixpunkt-Konvergenzrate.
\item Volumen-Schritt: $\lambda_{\mathrm{sphere}}=1/(1+2\pi)$; $\Delta(M)$ bleibt
  nach isotroper Skalierung null, wenn es vorher null war.
\end{itemize}

\subsection{Numerische Verifikation}

Die vier Kern-Lemmata werden von \texttt{run\_core\_lemma\_tests()} in
\texttt{eabc\_icosahedron\_test.py} geprüft (Toleranz $10^{-6}$):

\begin{center}
\begin{tabular}{clc}
\hline
Lemma & Testschlüssel & Status \\
\hline
1 (Isotropie-Lemma)        & \texttt{isotropy}     & PASS \\
2 (Projektor-Idempotenz)   & \texttt{idempotence}  & PASS \\
3 (Dualitätsstabilität)    & \texttt{duality}      & PASS \\
4 (Volumen-Kompatibilität) & \texttt{volume}       & PASS \\
\hline
\multicolumn{2}{r}{Gesamt:} & 4/4 PASS \\
\end{tabular}
\end{center}

%=============================================================================
\section*{Synthese und Schlussfolgerung}
\label{sec:conclusion}
%=============================================================================

\begin{center}
\label{box:conclusion}
\fbox{%
  \fbox{%
    \begin{minipage}{0.92\linewidth}
      \vspace{2mm}
      \begin{center}
        \textbf{\large Synthese der EABC-Renormierung}
      \end{center}
      \vspace{3mm}
      \begin{equation*}
        \textbf{Quantisierter Primdefekt}\;(\text{effektive Ebene }M_{\mathrm{eff}})\!:
        \quad
        p \;\longmapsto\;
        w_p \;=\;
        \Delta\!\bigl(24I_3 + w_p\,vv^{\mathsf T}\bigr),
        \qquad \|v\|=1.
      \end{equation*}
      \begin{equation*}
        \textbf{Globale Tensor-Restauration}\;(\text{Peano-Schalenstapel})\!:
        \quad
        \forall i<n:\;
        M_{\mathrm{eff}}\!\bigl(\mathcal{R}^*(\mathcal{S}^{(n)}_{\mathrm{def}})\bigr)_i
        \;=\;
        24I_3.
      \end{equation*}
      \vspace{2mm}
      \hrule
      \vspace{2mm}
      \small
      \textbf{Kernaussage.}
      Eine EABC-Primzahl $p>3$ mit $p\bmod 12\in\{1,5,7,11\}$ induziert auf einer
      Normschale eine quantisierte, messbare Rang-1-Anisotropie
      $w_p\in\{1,5,7,11\}$.
      Der komponentenweise Retraktions-Schnitt $\mathcal{R}^*$ führt das
      Schalensystem algebraisch stabil auf den isotropen Ikosaeder-Fixpunkt zurück.
      Die Satzkette ist in \texttt{eabc-renorm} maschinell verifiziert
      (Lean~4.29, \texttt{lake build} grün, \texttt{sorry}\,$=\,0$;
      Tabelle~\ref{tab:lean-theorems}).
      Metrische Skalierungen $r_n$ (Phase~B) sind interpretativ und berühren die
      Verifikation nicht.
      \smallskip

      \textit{Intellektuelle Linie:}
      Von Eulers polyedrischer Symmetrie über Riemanns Spektralgeometrie und
      Diracs gestörtem Vakuum zu 't~Hoofts projektiver Renormierungslogik ---
      hier als diskrete, maschinengeprüfte Satzkette auf dem Ikosaeder-Fixpunkt.
      \vspace{2mm}
    \end{minipage}%
  }%
}
\end{center}

%=============================================================================
\appendix
\section{Explizite Erzeuger-Verifikation}
\label{app:generators}
%=============================================================================

\subsection{Ikosaeder-Ecken und geometrische EABC-Zuweisung}

Mit $\varphi=(1+\sqrt5)/2$ seien die zwölf normierten Ikosaeder-Ecken
\[
  v_i \in
  \Bigl\{
  \frac{(0,\pm1,\pm\varphi)}{\|(0,1,\varphi)\|},
  \frac{(\pm1,\pm\varphi,0)}{\|(1,\varphi,0)\|},
  \frac{(\pm\varphi,0,\pm1)}{\|(\varphi,0,1)\|}
  \Bigr\}.
\]
Die Koordinaten-Restklasse $r(v)$ mit Skala $s=6$ teilt die Ecken in vier
Bahnen zu je drei Elementen mit Restklassen $2,4,8,10\bmod 12$.
Die Abbildung \emph{nächste EABC-Restklasse} liefert
\[
  2\mapsto E\;(w=1),\quad
  4\mapsto A\;(w=5),\quad
  8\mapsto B\;(w=7),\quad
  10\mapsto C\;(w=11).
\]

\subsection{Direkte Berechnung von $M$}

Am geometrischen Fixpunkt $I_{12}^{\mathrm{EABC}}$ gilt
\[
  M
  \;=\;
  \sum_{i=1}^{12} w_{\sigma(i)}\, v_i v_i^{\mathsf T}
  \;\in\;
  \Q(\sqrt5)^{3\times 3}.
\]
Eine direkte Summation über die zwölf äußeren Produkte (symbolisch in
$\varphi$) ergibt exakt
\[
  M = 24\,I_3.
\]
Damit ist $\Delta(M)=0$ und $\lambda_1=\lambda_2=\lambda_3=24$ unmittelbar
sichtbar, unabhängig von der Symmetrie-Rechnung im Beweis von Lemma~1.

\subsection{Erzeuger $R_3$ und $R_5$}

Die $2\pi/3$-Drehung um den Flächenschwerpunkt $c=\tfrac{1}{\sqrt3}(1,1,1)$
hat die Matrix
\[
  R_3
  \;=\;
  \begin{pmatrix}
    0&1&0\\ 0&0&1\\ 1&0&0
  \end{pmatrix}
  \;\in\; SO(3).
\]
Die $2\pi/5$-Drehung um die Ecke $v_0$ (Rodrigues-Formel) liefert
\[
  R_5
  \;=\;
  \begin{pmatrix}
    \tfrac{-1+\varphi}{2} & \tfrac{1+\varphi}{2} & -\tfrac12 \\[4pt]
    -\tfrac{1+\varphi}{2} & \tfrac12 & \tfrac{-1+\varphi}{2} \\[4pt]
    \tfrac12 & \tfrac{-1+\varphi}{2} & \tfrac{1+\varphi}{2}
  \end{pmatrix}
  \;\in\; SO(3),
\]
wobei man $\cos(2\pi/5)=\tfrac{-1+\varphi}{2}$ und
$\sin(2\pi/5)=\tfrac{\sqrt{10+2\sqrt5}}{4}$ verwendet.

\subsection{Verifikation der $\mathcal{I}$-Invarianz}

Da $M=24I_3$ exakt gilt, folgt für jedes $R\in SO(3)$ sofort
\[
  R\,M\,R^{\mathsf T}
  \;=\;
  24\,R\,R^{\mathsf T}
  \;=\;
  24\,I_3
  \;=\;
  M.
\]
Insbesondere für die Erzeuger $R_3,R_5\in A_5$:
\[
  R_k\,M\,R_k^{\mathsf T} = M,
  \qquad k=3,5.
\]
Damit ist Schritt~(iii) im Beweis des Isotropie-Lemmas vollständig explizit
nachgerechnet; die Untergruppe $\mathcal{G}\subset\mathcal{I}$ aus dem
Hauptbeweis enthält $R_3,R_5$ und daher ganz $\mathcal{I}\cong A_5$.

\end{document}
